% sec-A083829.tex -- k = 3 section of the unified kp+1 paper.
% Standalone fragment, \input by A083833-UNIFIE-paper.tex.
% Environments and notation inherited from sec-carry.tex; the
% guards below are no-ops once defined there.
\providecommand{\IsPal}{\mathrm{IsPal}}
\makeatletter
\@ifundefined{theorem}{\newtheorem{theorem}{Theorem}}{}
\@ifundefined{corollary}{\newtheorem{corollary}[theorem]{Corollary}}{}
\@ifundefined{remark}{\newtheorem{remark}[theorem]{Remark}}{}
\makeatother

\section{Multiplier $k = 3$: A083829 and its image A083830}
\label{sec:k3}

On the entry A083829~\cite{oeis} -- palindromes $p$ such
that $q = 3p+1$ is again a palindrome -- Robert Israel
(February 23, 2023, tagged ``It appears'') recorded four
observations, numbered here P1--P4 and quoted in
Section~\ref{sec:intro}. The carry chain of
Section~\ref{sec:carry} at $k = 3$ has
$c_j = \bigl\lfloor\bigl(3\,(p \bmod 10^{j}) + 1\bigr)
/10^{j}\bigr\rfloor$, so $c_0 = 1$ and $c_j \le 2$.

\begin{theorem}[Israel's observations on A083829]
\label{thm:k3-israel}
Let $\IsPal(p,m)$ and $\IsPal(3p+1,m')$ with $m \ge 1$.
Then:
\begin{itemize}
\item[(P1)] if $m$ is even, every digit of $p$ is $7$
($p$ has the shape $77\ldots7$ of A002281);
\item[(P2)] if $m \ge 2$ and some digit of $p$ differs
from $7$ -- equivalently $p$ is not the all-$7$
repdigit -- the leading pair $10\,a_{m-1} + a_{m-2}$
lies in $\{14, 15, 16, 24, 25, 26\}$;
\item[(P3)] no digit of $p$ lies in $\{3, 8, 9\}$;
\item[(P4)] if some digit of $p$ equals $7$, then every
digit equals $7$.
\end{itemize}
\end{theorem}

\begin{corollary}[alphabet of the image; A083830]
\label{cor:k3-alphabet}
Under the same hypotheses, every digit of $q = 3p+1$
avoids $\{0, 6, 9\}$, i.e.\ lies in
$\{1,2,3,4,5,7,8\}$.
\end{corollary}

\begin{proof}
Two branches of the shared carry lemma. \emph{Same
length} ($m' = m$): Lemma~\ref{lem:same} forces a
$2$-periodic chain, $m$ odd, $a_0 \in \{1,2\}$ and
$c_1 = 0$, confining the digits to
$a_j \in \{0,1,2\}$ at even positions and
$a_j \in \{4,5,6\}$ at odd positions. \emph{Overflow}
($m' = m+1$): the head equation forces $a_0 = 7$ and
$c_1 = 2$; inductively the joint constraints admit one
spurious state $(a_{j+1}, c_{j+2}) = (3,1)$, locally
consistent but a dead end -- at the next position the
shifted mirror congruence forces the digit $6$, and it
would require $0 = \lfloor(3 \cdot 6 + c)/10\rfloor =
\lfloor(18 + c)/10\rfloor$ -- so only
$(7,2)$ survives in the interior: every digit is $7$,
$p = 77\ldots7$ and $q = 23\ldots32$. Now P1 follows
because even $m$ rules out the same-length case, which
needs $m$ odd; P3 from the alphabet union
$\{0,1,2\} \cup \{4,5,6\} \cup \{7\}$; P4 because a digit
$7$ excludes the same-length alphabet; and P2 because in
the non-repdigit same-length case $m$ is odd, so position
$m-2$ is odd: $a_{m-1} = a_0 \in \{1,2\}$ and
$a_{m-2} \in \{4,5,6\}$. For the image alphabet:
same-length case $b_j = 3a_j + 1 \in \{1,4,7\}$ at even
$j$ and $b_j = 3a_j \in \{12,15,18\}$ at odd $j$, digits
in $\{1,2,4,5,7,8\}$; overflow case $q = 23\ldots32$,
digits in $\{2,3\}$.
\end{proof}

\begin{remark}[P1 and the ``Includes A002281'' clause]
\label{rem:k3-iff}
P1 states the necessity direction; in the kernel the
even-length classification is the literal iff
\emph{even-length member $\Leftrightarrow$ repdigit
$77\ldots7$}, so Israel's ``Includes A002281'' clause --
every repdigit $77\ldots7$ is a member, image
$23\ldots32$ -- is itself certified
(Table~\ref{tab:lean}). P1--P4 are universal properties
of the members, not an enumeration of the sequence. The
$37$ published terms agree: their even-length members
are exactly $77$, $7777$, $777777$.
\end{remark}
